\section{背景}
\label{sec:background}

\paragraph{告警分诊} 告警分诊针对检测产生的每条告警，判断实际发生了什么、是否需要响应~\citep{nodoze}。它与攻击调查不同：攻击调查从一个关注点出发重建完整的攻击过程~\citep{kairos,clouseau}，而分诊以单一结论结束，即告警的原因及支持它的证据。分析员通过查询环境得出结论，而查询在成本和区分候选解释的能力上各不相同：有些查询便宜且对某个解释具有决定性，另一些（例如威胁情报查询）既慢，又对几种攻击返回同样的答案。设想一条告警报告 Web 层认证失败和 HTTP 错误激增。它看起来像凭证填充，但同样的症状也可能来自 SQL 注入探测、授权的漏洞扫描、被误判的健康检查，或者暴露在公网上的管理面板；读取管理面板的配置只需一次便宜的查询，就能立刻判定最后一种解释，但只有始终把这个解释放在视野里的调查者才会运行它。因此，查询的顺序决定调查多快结束，调查者还必须判断证据何时已经足够。

\paragraph{本地 LLM 智能体} LLM 智能体把语言模型与工具结合起来：每一步，模型读取任务和已有观测，然后调用一个工具或结束任务~\citep{react}。Qwen2.5 和 Llama~3.1 等开源权重模型可以在组织自己的硬件上运行，4 比特权重下单块 GPU 最多可运行 72B 参数，这让遥测留在本地，却让循环的每一步都依赖一个较弱的模型。

\paragraph{对智能体输入的注入} 读取检索内容的智能体会被其中隐藏的指令操纵，即间接提示注入~\citep{greshake2023}，AgentDojo~\citep{agentdojo} 在工具型智能体上对此进行了基准测试。日志注入通过未转义的用户可控字段伪造或破坏日志记录，是同一问题更早的实例~\citep{owasp_loginjection}。在分诊中两者都适用，因为请求载荷、用户代理、命令行和工单文本都由引发告警的一方书写，并被复制进智能体读取的日志。
